\chapter*{The Empty Tomb}

Shortly after the Crucifixion, another rather mysterious pair of figures came onto the scene. One of them was Joseph of Arimathea, ``a prominent council member, who was himself waiting for the kingdom of God'' (Mk 15:43); in the apocryphal Gospel of Peter, he is characterized as a ``friend of Pilate and of the Lord.'' It was this man who received the procurator's permission to collect the body of an executed criminal---the ``King of the Jews''---and bury him in his own tomb. The Evangelists depict Joseph as a ``secret disciple of Christ'' (Jn 19:38), though nowhere in the text of the New Testament are there any previous indications of him having any contact with Jesus or the Apostles---either before or, strangest of all, after the burial. Granted, there is one legend that claims he was the first to preach the Gospel in Britain, but even the official Church considers this unlikely.\footnote{``Joseph,'' \emph{Bible Encyclopedia}, vol. 2, pp. 111-115.}

Assisting him was Nicodemus, another member of the Sanhedrin. Unlike Joseph, Nicodemus had actually met Christ twice before; on the first occasion, he spent a whole night conversing with him (Jn 3:1-21), and in another episode, he openly came out in his support before other members of the Sanhedrin (Jn 7:50-52). Thus, we can assume that the initiator of the burial was not actually Joseph, but Nicodemus. The legend of Nicodemus' eventual fate seems far more modest in comparison (``he was subsequently baptized by the Apostles''\footnote{Nicodemus,'' \emph{Bible Encyclopedia}, vol. 3, p. 71.}), but for that reason, it sounds far more believable.

The fact remains, however, that two prominent members of the local political establishment rather openly defied the Judean authorities, and did so at the very moment it was clear playtime was over; meanwhile, all of Christ's ``official'' disciples were quaking in their boots and concerned solely with their own survival. So what prompted them to do this? Were they merely ``waiting for the Kingdom of God?'' Hmm... For what it's worth, Joseph and Nicodemus actually had something to lose here---unlike the disciples, who had the social status of bums.

But all of this is just a preamble. Now it's time for us to move straight to one of the key events in our story: the disappearance of Christ's body from a guarded tomb. This tomb, which (as mentioned above) belonged to Joseph of Arimathea, was a new burial vault that had just been carved out from the rock, and was located in a fairly isolated place close to Jerusalem; the latter detail made the guards' task much easier. The entrance to the tomb, which was covered by a heavy boulder (McDowell gives its weight as 1.5 tons) and secured with an imperial seal, was guarded by Roman soldiers. The fact they were specifically Roman (that is, highly disciplined, and neutral with regard to intra-Judean disputes) plays a significant role in McDowell and Gladkov's arguments, and they give detailed and convincing justifications for it.\footnote{Here we could make the following objection: in response to the high priests' request, ``Pilate says to them, `You have a guard; go your way, make it as secure as you know how''' (Mt 27:65). In principle, the procurator's words could be interpreted as a refusal (``You guys have your own temple guard---go deal with your own scandals!''). McDowell, however, makes a purely linguistic argument that in the original Gospel text (unlike many of its translations into European languages), all three verbs should be read as being in the same mood: the imperative. Therefore, what this sentence really sounds like is ``have a guard; go your way, make it as secure as you know how,'' that is: ``Take some soldiers and act as you see fit going forward.'' (K.E.)}

At dawn on the third day, the guards discover to their astonishment that the stone has been rolled aside and the tomb is now empty, and immediately inform the high priests of this:

\begin{smallquote}
    When they had assembled with the elders and consulted together, they gave a large sum of money to the soldiers, saying, ``Tell them, `His disciples came at night and stole Him away while we slept.' And if this comes to the governor's ears, we will appease him and make you secure.'' So they took the money and did as they were instructed; and this saying is commonly reported among the Jews until this day. (Mt 28:12-15)
\end{smallquote}
Pilate subsequently figured that the high priests' bribe proved the soldiers' innocence, and decided not to punish the latter.

I am perfectly willing to admit that there were some members of the Judean hierarchy who genuinely believed in the divine, messianic essence of Christ. The Romans, however, are a different story altogether; everything historical and literary sources tell us about their mentality suggests that Roman society at that time was actually quite secular in nature. It is precisely this fact that is typically used to explain ancient civilization's fundamental inability to counteract the moral expansion of Christianity several decades later. For this reason, the very thought that the pragmatic Romans, who failed to show piety even to their own gods, were seriously willing to accept the Jewish legends of a Messiah who would rise again three days after his death---the very thought of that seems absurd to me. That's point number one.

Point number two: it bears mentioning that the disciplinary code of the Roman army was quite severe. McDowell notes that soldiers who had fallen asleep at their post were subject to execution without exception (regardless of the result of their negligence). With both these points in mind, we have no choice but to admit that \emph{every} action taken by \emph{every} party involved in this incident---the guards, the high priests, and Pilate---is completely absurd (no matter the actual reason for the tomb being empty). It's like they each agreed to do the one thing that would work most effectively against their own interests. Judge for yourself.

For example, take how the guards discover the loss of the ``protected object'' in the morning. If the soldiers had really kept watch all night, they could be sure this appalling circumstance wasn't the result of any complicated games on the part of the Judean authorities. On the other hand, it was the guards themselves who discovered the loss, and not one of the officers. In such a situation, the most logical thing to do would be to change shifts at dawn as planned and brazenly report that everything was in order---maybe they'd actually buy it. If the Judean authorities go on to find the open tomb and start making a fuss, just put on a poker face and declare that when the guards changed shifts---cross my heart!---everything was A-OK, and I have no way of knowing what went on there afterwards; those Asian ragheads can go look after their own stiffs, for all I care---I've had it up to here with them, Your Honor! 

But still, let's suppose the commander of the guard isn't a savvy sergeant major of Soviet vintage, but what one might call ``an honest, yet stupid young lady.'' Why then, however, does this disciplined Roman bring his penitential report not to his commanding officers, as regulations require, but to the local authorities? Could he really be counting on the protection of Pilate's hated Sanhedrin? Well, in that case, he's just plain dumb. Then things get even stupider: the Sanhedrin offers to pay the guards a modest sum to sign their own death sentence (admitting they fell asleep at their post). The guards agree and go about earning their bribe honestly: they shout from the rooftops about an act of official misconduct they have supposedly committed---on the off chance that Roman leadership is still woefully ignorant of their exploits.

As for the high priests: if they really wanted to dispute the fact of the resurrection, they should have wasted no time in accusing the soldiers of sleeping at their post and missing the theft of the body. Without a doubt, the Judeans knew that in the eyes of Roman officers, this incident simply couldn't have any other explanation, and the soldiers' babble about ``miracles'' would only further prove their guilt. Just try to imagine a modern-day general whose underlings attempt to foist a report like this onto him: ``While we were standing watch, a flying saucer landed next to our post and paralyzed the personnel with a blue ray of light, at which point little green men entered the depot we were guarding and carried out 42 assault rifles and 12 boxes of hand grenades.'' Three guesses how the general would react? In the very act of negotiating with and bribing the guards, however, the high priests all but admitted to their belief in Christ's resurrection---as Pilate would later acknowledge himself.

And now, the procurator. First, his subordinates have fallen asleep at their post and lost track of the ``protected object,'' an act deserving of the death penalty. Even worse, they're taking bribes from the local poobah and carrying out his orders; in any army in the world, this would be considered an even worse crime than the original offense. In our case, however, two wrongs surprisingly make a right: the guards are ultimately spared any punishment whatsoever by Pilate---he is apparently satisfied by their account of the body's miraculous disappearance. I can't help but think back to another episode in this regard. Later on, the Apostle Peter just as mysteriously disappears from a guarded prison cell; however, Herod never even thinks of taking the ``miraculous'' nature of this event into account, and immediately executes the guards---in accordance with regulations (Ac 12:19).

Speaking of St. Peter's disappearance:

\begin{smallquote}
    That night Peter was sleeping, bound with two chains between two soldiers; and the guards before the door were keeping the prison. [...] An angel [...] struck Peter on the side and raised him up, saying, ``Arise quickly!'' And his chains fell off his hands. Then the angel said to him, ``Gird yourself and tie on your sandals. [...] Put on your garment and follow me.'' [...] When they were past the first and the second guard posts, they came to the iron gate that leads to the city, which opened to them of its own accord; and they went out and went down one street, and immediately the angel departed from him. (Ac 12:6-10)
\end{smallquote}
I don't know about you, but I'm personally on Herod's side here: while of course, these events can be considered \emph{mysterious}, by no means can they be considered \emph{miraculous}. And if Herod was interested in the identity of the Angel behind this caper (and you know he was!), then he certainly didn't consult his theologian on staff, but the chief of his security service. It should be noted that this wasn't the first time powerful messengers of a Higher Power pulled tricks like this to set prisoners free (for example, Ac 5:18-24). It's no surprise that the patience of the Tetrarch of Galilee \emph{und} Perea wore thin (``Am I king or am I not king?''\footnote{A quote from \emph{Tsar Fyodor Ioannovich} (see note 47).}), and he tried putting these utterly insolent Angels in their place.

In refuting the theory that Christ's body was actually stolen by his disciples, McDowell, among other things, brings up the following point. He argues that it was unthinkable for a citizen of Judea to break the Roman seal that was placed on the tomb. The penalty for this under Roman law was inverted crucifixion, and the Empire's secret services knew neither sleep nor rest until the perpetrator was captured. While this information is certainly interesting, McDowell strangely fails to observe that, like it or not, the seal was still broken, and as it happened, there was no investigation into it---not even the most perfunctory one. That being said, it would be easier to quickly crucify a couple of disciples and close the case of the broken seal that way; this was exactly how Nero would close the case of Rome's ``burning'' sometime later. In reality, as we know, nothing of the sort ever came to pass, and this fits in perfectly with the general picture of the Roman authorities' strange benevolence towards Christ and his followers.

I can only find one explanation for this strange turn of events. The testimony of the negligent guards, as well as some other details of the incident---were they to surface in the course of an official investigation---would for some reason be so uncomfortable for Pilate that he preferred to let the entire affair blow over. In this episode in general, it seems as if both the Roman and the Judean authorities are acting in accordance with a tacit agreement, simultaneously attempting to avert a scandal that might be fraught with dangerous revelations. We could assume, for example, that both of these ``frenemies'' were the object of serious blackmail on the part of a certain ``third force.''

In this regard, the apocryphal Gospel of Peter is quite remarkable. It was apparently written with the aim of dispelling all doubt as to the fact of the Resurrection; here, it occurs directly in front of a large group of people, and with the direct participation of two angels:

\begin{smallquote}
    The heavens were opened and that two males who had much radiance had come down from there and come near the sepulcher. But that stone which had been thrust against the door, having rolled by itself, went a distance off the side; and the sepulcher opened, and both the young men entered. [...] Again [the guards] see three males who have come out from the sepulcher, with the two supporting the other one, and a cross following them, and the head of the two reaching unto heaven, but that of the one being led out by a hand by them going beyond the heavens.\footnote{Raymond E. Brown, \emph{The Death of the Messiah: from Gethsemane to the Grave}, v. 2, p. 1320}
\end{smallquote}
In the end, the Evangelist overdid it a little---the number of miracles per square meter of text fails to abide by even the most lenient standards of realism. Perhaps this is why this version was considered apocryphal.

However, it contains one particularly curious detail: this is the only Gospel out of all of them where the guard is combined. It consists of Roman legionnaires led by the centurion Petronius\footnote{The Greek word \emph{kentyrion} (``centurion''), which is used in the text of this Gospel, forces us to assume that its author was a Roman, and not a resident of Palestine. The issue here is that in all the canonical Gospels, the Roman officer ranks of centurion and military tribune are rendered as, respectively, \emph{hekatontarch} (for example, Lk 7:2-9) and \emph{chiliarch} (for example, Jn 18:12). (K.E.)} as well as Jewish elders and scribes, who jointly keep watch---a stretch if there ever was one!---from inside a single tent. It's quite telling that the unknown Evangelist made sure (among other things) to include such a detail when naively composing his ``extra-convincing'' version of the Resurrection. Clearly, a combined guard (which would ensure that each faction kept the other in check) would put people's doubts to rest far more convincingly than the exclusively Roman guard in the canonical version.

Still, let's take another look at the series of events that led to the mysterious emptying of the tomb, starting from the very beginning. Let's rewind the tape a little to Friday afternoon, when Nicodemus and Joseph collect Christ's body and perform all the rituals over it that are required by Jewish custom. McDowell himself gives a detailed description of such rituals, which I read with great (purely ethnographical) interest. Jews would wrap their dead in numerous layers of fabric strips, which were soaked in aromatic substances. Moreover, the fabric could hold up to 40-50 kilograms' worth of resinous material; the corpse would eventually be enclosed in a thick shell of fabric and resin. The ``empty burial shrouds'' that Peter and John discovered in Christ's tomb were something like an empty cocoon a butterfly had fluttered out of.

Before nightfall, Nicodemus and Joseph manage to transport the body to a tomb outside the city and seal the entrance with a stone; all of this occurs in the presence of Jesus' female companions:

\begin{smallquote}
    And the women who had come with Him from Galilee followed after, and they observed the tomb and how His body was laid. (Lk 23:55)
\end{smallquote}
Speaking of the stone: you have to think that it wasn't anywhere near as heavy as people sometimes claim (Mk 16:4)---after all, it was quietly rolled into place that evening by two men who clearly weren't weightlifters. The next day (Saturday), the Judean leadership comes to a realization:

\begin{smallquote}
    The chief priests and Pharisees gathered together to Pilate, saying, ``Sir, we remember, while He was still alive, how that deceiver said, `After three days I will rise.' Therefore command that the tomb be made secure until the third day, lest His disciples come by night and steal Him away, and say to the people, `He has risen from the dead.''' (Mt 27:62-64)
\end{smallquote}
For this occasion, the Jewish high priests decide to desecrate their sacred day of
Sabbath rest; with soldiers in tow, they head out to the tomb and place a seal on the
stone at its entrance. The Roman guard stationed at the tomb...

Wait!!! Hold on, run that by me one more time!! ``The next morning, the high priests...'' There it is: that means that \emph{from early Friday evening to late Saturday morning, the unsealed tomb was left completely unsupervised}. It makes me wonder: just what did the Roman soldiers end up ``taking into safekeeping?''